% Two masses: CHM on all balls; Gauss not killed when both sit inside.
\documentclass[11pt]{article}
\usepackage[margin=1.05in]{geometry}
\usepackage{amsmath,amssymb,amsthm}
\usepackage[colorlinks=true,linkcolor=blue,citecolor=blue,urlcolor=blue]{hyperref}

\newtheorem*{admission}{Admission}

\title{\textbf{Two Masses: CHM Wins on All Balls.\\
Gauss Survives When Both Sit Inside.}}
\author{Robert W.\ McGwier\\
\small CTO, Cohere Technology Group\\
\small GitHub: \texttt{n4hy}}
\date{15 August 2026 --- Version 1}

\begin{document}
\maketitle

\begin{abstract}
\noindent
An opposite-energy dipole cannot be built on the Fermi-sea
vacuum. Two positive occupation transfers are two masses.

$N=12$, $512$ balls. CHM beats flat ($R=0.044$ vs $0.159$).
On the $11$ balls that contain both sources, enclosed
energy is nearly constant and the CHM gap collapses
($0.107$ vs $0.115$).

Auditor, $N=10$: all-ball C2 CONFIRMED. Both-inside C2b
REFUTED (flat wins).

The all-ball CHM win is the pair entering or leaving the
ball. It is not a boost-versus-Gauss theorem. Not $8\pi G$.
Not Einstein. Not de~Sitter.
\end{abstract}

\begin{admission}
I wanted a dipole. The vacuum forbids it. I ran two masses
instead, and I did not cash eleven barely-separated points
as linearized gravity.
\end{admission}

\end{document}
